\chapter{Metody sumowania}
\label{ma1:ch:summability}
Zwykła zbieżność szeregu oznacza zbieżność ciągu jego sum częściowych. Metody sumowania przypisują wartości niektórym szeregom rozbieżnym, zastępując tę granicę inną procedurą. Metodę nazywamy regularną, jeżeli każdemu zbieżnemu szeregowi przypisuje jego zwykłą sumę. W tym rozdziale indeksujemy ciągi od zera.

\section{Metoda Cesàro}
\begin{definicja}
\label{ma1:def:cesaro}
Niech $K=\R$ lub $K=\C$, $(a_n)_{n\ge0}$ będzie ciągiem w $K$, a $s_n=\sum_{j=0}^n a_j$. Szereg $\sum_{n\ge0}a_n$ nazywamy sumowalnym metodą Cesàro (rzędu pierwszego) do $S\in K$, jeżeli
\[
 c_n=\frac1{n+1}\sum_{k=0}^n s_k\longrightarrow S.
\]
Liczbę $S$ nazywamy sumą Cesàro szeregu.
\end{definicja}
Zmieniając kolejność sumowania w skończonej sumie, otrzymujemy $c_n=\sum_{j=0}^n(1-j/(n+1))a_j$, ponieważ $a_j$ występuje w dokładnie $n-j+1$ sumach częściowych.

\begin{lemat}
\label{ma1:lem:cesaro-means}
Jeżeli ciąg $(x_n)_{n\ge0}$ liczb rzeczywistych lub zespolonych dąży do $L$, to $(n+1)^{-1}\sum_{k=0}^n x_k\to L$.
\end{lemat}
\begin{proof}
Dla $\varepsilon>0$ wybieramy $m\ge1$ takie, że $|x_k-L|<\varepsilon/2$ dla $k\ge m$. Oznaczmy $A=\sum_{k=0}^{m-1}|x_k-L|$. Dla $n\ge m$ mamy
\[
 \left|\frac1{n+1}\sum_{k=0}^n x_k-L\right|
 \le\frac{A}{n+1}+\frac{n-m+1}{n+1}\frac{\varepsilon}{2}.
\]
Drugi składnik nie przekracza $\varepsilon/2$, a pierwszy jest mniejszy od $\varepsilon/2$ dla dostatecznie dużych $n$.
\end{proof}
\begin{wniosek}
\label{ma1:cor:cesaro-regular}
Każdy zbieżny szereg liczb rzeczywistych lub zespolonych jest sumowalny metodą Cesàro do swojej zwykłej sumy.
\end{wniosek}
\begin{proof}
Stosujemy lemat \ref{ma1:lem:cesaro-means} do ciągu sum częściowych szeregu.
\end{proof}
\begin{przyklad}
\label{ma1:ex:grandi-cesaro}
Dla $a_n=(-1)^n$ mamy $s_{2m}=1$, $s_{2m+1}=0$ oraz
$c_{2m}=(m+1)/(2m+1)$, $c_{2m+1}=1/2$. Oba podciągi średnich dążą do $1/2$, więc cały ciąg średnich dąży do $1/2$: dla każdej dokładności wystarczy wziąć wspólny próg dla obu podciągów. Szereg nie jest zbieżny, lecz jego suma Cesàro wynosi $1/2$.
\end{przyklad}
\begin{przyklad}
\label{ma1:ex:cesaro-fails}
Dla $a_n=1$ mamy $s_n=n+1$ oraz $c_n=(n+2)/2\to+\infty$. Szereg nie jest sumowalny metodą Cesàro.
\end{przyklad}
\begin{twierdzenie}
\label{ma1:thm:cesaro-term}
Jeżeli szereg $\sum_{n\ge0}a_n$ liczb rzeczywistych lub zespolonych jest sumowalny metodą Cesàro, to $a_n/n\to0$ dla $n\ge1$.
\end{twierdzenie}
\begin{proof}
Dla średnich $c_n$ i sum częściowych $s_n$ mamy $s_n=(n+1)c_n-nc_{n-1}$ dla $n\ge1$. Jeżeli $c_n\to S$, to $s_n/n=(1+1/n)c_n-c_{n-1}\to0$. Dla $n\ge2$ mamy $a_n/n=s_n/n-((n-1)/n)(s_{n-1}/(n-1))\to0$.
\end{proof}
Warunek ten nie jest wystarczający, jak pokazuje przykład \ref{ma1:ex:cesaro-fails}.

\section{Metoda Abela}
\begin{definicja}
\label{ma1:def:abel-summation}
Niech $(a_n)_{n\ge0}$ będzie ciągiem liczb rzeczywistych lub zespolonych. Szereg $\sum_{n\ge0}a_n$ jest sumowalny metodą Abela do $S$, jeżeli dla każdego $r\in[0,1)$ szereg $\sum_{n\ge0}a_nr^n$ jest zbieżny oraz
\[
 \lim_{r\to1^-}\sum_{n=0}^\infty a_nr^n=S.
\]
Przy $r=0$ przyjmujemy $r^0=1$.
\end{definicja}
Najpierw obliczamy sumę szeregu dla ustalonego $r<1$, dopiero później przechodzimy do granicy.
\begin{lemat}
\label{ma1:lem:polynomial-geometric}
Dla każdego $r\in[0,1)$ szeregi $\sum_{n\ge0}r^n$ i $\sum_{n\ge0}(n+1)r^n$ są zbieżne i mają odpowiednio sumy $1/(1-r)$ i $1/(1-r)^2$.
\end{lemat}
\begin{proof}
Pierwszy wzór wynika ze skończonej sumy geometrycznej i granicy $r^{n+1}\to0$. Drugi wynika z twierdzenia o iloczynie Cauchy'ego dwóch bezwzględnie zbieżnych szeregów geometrycznych: współczynnik o indeksie $n$ jest sumą $n+1$ składników równych $r^n$.
\end{proof}
\begin{lemat}
\label{ma1:lem:abel-weights}
Niech $(x_n)_{n\ge0}$ będzie ciągiem liczb rzeczywistych lub zespolonych zbieżnym do $L$. Dla $r\in[0,1)$ wybierzmy jedną z rodzin wag
$w_n(r)=(1-r)r^n$ lub $w_n(r)=(1-r)^2(n+1)r^n$.
Wówczas szereg $\sum_{n\ge0}w_n(r)x_n$ jest bezwzględnie zbieżny oraz jego suma dąży do $L$ przy $r\to1^-$.
\end{lemat}
\begin{proof}
Wagi są nieujemne, ich suma wynosi $1$ na mocy lematu \ref{ma1:lem:polynomial-geometric}, a dla ustalonego $n$ zachodzi $w_n(r)\to0$ przy $r\to1^-$. Ciąg $(x_n)$ jest ograniczony, co przez porównanie daje bezwzględną zbieżność. Dla $\varepsilon>0$ wybieramy $m\ge1$ tak, by $|x_n-L|<\varepsilon/2$ dla $n\ge m$. Wówczas
\[
 \left|\sum_{n=0}^\infty w_n(r)x_n-L\right|
 \le\sum_{n=0}^{m-1}w_n(r)|x_n-L|+\frac{\varepsilon}{2}.
\]
Skończona suma po prawej dąży do zera. Otrzymujemy tezę.
\end{proof}
\begin{twierdzenie}
\label{ma1:thm:abel-regular}
Każdy zbieżny szereg liczb rzeczywistych lub zespolonych jest sumowalny metodą Abela do swojej zwykłej sumy.
\end{twierdzenie}
\begin{proof}
Niech $s_n=\sum_{j=0}^n a_j\to S$, $s_{-1}=0$. Ciąg $(s_n)$ jest ograniczony, a więc $(a_n)$ również, ponieważ $a_n=s_n-s_{n-1}$. Dla $0\le r<1$ szeregi $\sum a_nr^n$ i $\sum s_nr^n$ są bezwzględnie zbieżne przez porównanie z szeregiem geometrycznym. Odejmując szeregi i przesuwając indeks, mamy
\[
 \sum_{n=0}^\infty a_nr^n=(1-r)\sum_{n=0}^\infty s_nr^n.
\]
Lemat \ref{ma1:lem:abel-weights} dla pierwszej rodziny wag daje granicę $S$.
\end{proof}
\begin{przyklad}
\label{ma1:ex:grandi-abel}
Dla $a_n=(-1)^n$ szereg $\sum_{n\ge0}(-r)^n$ jest bezwzględnie zbieżny dla $0\le r<1$ i ma sumę $1/(1+r)$, która dąży do $1/2$. Dla $a_n=1$ suma $1/(1-r)$ dąży do $+\infty$, więc skończona suma Abela nie istnieje.
\end{przyklad}

\section{Porównanie metod}
\begin{twierdzenie}
\label{ma1:thm:cesaro-implies-abel}
Jeżeli szereg $\sum_{n\ge0}a_n$ liczb rzeczywistych lub zespolonych jest sumowalny metodą Cesàro do $S$, to jest sumowalny metodą Abela do $S$.
\end{twierdzenie}
\begin{proof}
Oznaczmy $s_n=\sum_{j=0}^n a_j$, $b_n=\sum_{j=0}^n s_j=(n+1)c_n$, $s_{-1}=b_{-1}=0$. Ponieważ $c_n\to S$, istnieje $M\ge0$ takie, że $|c_n|\le M$ dla wszystkich $n$. Stąd $|b_n|\le M(n+1)$ oraz
$|s_n|=|b_n-b_{n-1}|\le M(2n+1)$ dla $n\ge0$. Z kolei $|a_n|=|s_n-s_{n-1}|\le4M(n+1)$, również dla $n=0$. Z lematu \ref{ma1:lem:polynomial-geometric} wynika bezwzględna zbieżność wszystkich trzech szeregów $\sum a_nr^n$, $\sum s_nr^n$, $\sum b_nr^n$ dla każdego $0\le r<1$. Możemy zatem dwukrotnie odjąć szereg przesunięty o jeden indeks:
\[
 \sum_{n=0}^\infty a_nr^n=(1-r)\sum_{n=0}^\infty s_nr^n
 =(1-r)^2\sum_{n=0}^\infty(n+1)c_nr^n.
\]
Lemat \ref{ma1:lem:abel-weights} dla drugiej rodziny wag daje granicę $S$.
\end{proof}
\begin{przyklad}
\label{ma1:ex:abel-not-cesaro}
Dla $a_n=(-1)^n(n+1)$ każda para $a_{2j}+a_{2j+1}$ wynosi $-1$. Stąd $s_{2m}=m+1$, $s_{2m+1}=-(m+1)$. Każda para kolejnych sum częściowych $s_{2j}+s_{2j+1}$ wynosi zero, więc
$c_{2m}=(m+1)/(2m+1)$, $c_{2m+1}=0$. Średnie mają podciągi o różnych granicach $1/2$ i $0$, zatem szereg nie jest sumowalny metodą Cesàro.
Dla $0\le r<1$ szereg $\sum a_nr^n$ jest bezwzględnie zbieżny na mocy lematu \ref{ma1:lem:polynomial-geometric}. Iloczyn Cauchy'ego dwóch szeregów $\sum_{n\ge0}(-r)^n$ daje
\[
 \sum_{n=0}^\infty(-1)^n(n+1)r^n=\frac1{(1+r)^2}\longrightarrow\frac14
 \quad(r\to1^-).
\]
Szereg jest sumowalny metodą Abela do $1/4$. Odwrotność twierdzenia \ref{ma1:thm:cesaro-implies-abel} jest fałszywa.
\end{przyklad}

\begin{center}
\begin{tikzpicture}[node distance=7mm,every node/.style={font=\small,align=center},>=Stealth]
\node[draw,rounded corners,inner sep=6pt] (a) {Zwykła zbieżność};
\node[draw,rounded corners,below=of a,inner sep=6pt] (b) {Sumowalność metodą Cesàro};
\node[draw,rounded corners,below=of b,inner sep=6pt] (c) {Sumowalność metodą Abela};
\draw[->] (a)--(b);
\draw[->] (b)--(c);
\node[right=8mm of b,text width=5cm] {Przykład \ref{ma1:ex:grandi-cesaro} wyklucza pierwszą odwrotność.};
\node[right=8mm of c,text width=5cm] {Przykład \ref{ma1:ex:abel-not-cesaro} wyklucza drugą odwrotność.};
\end{tikzpicture}
\end{center}

\section{Liniowość i usuwanie skończonego początku}
\begin{twierdzenie}
\label{ma1:thm:summability-linear-shift}
Niech $K=\R$ lub $K=\C$. Ustalmy jedną z metod: Cesàro rzędu pierwszego lub Abela.
\begin{enumerate}[label=(\roman*)]
\item Jeżeli szeregi $\sum_{n\ge0}a_n$ i $\sum_{n\ge0}d_n$ w $K$ są sumowalne tą metodą do $A$ i $D$, to dla dowolnych $\alpha,\beta\in K$ szereg $\sum_{n\ge0}(\alpha a_n+\beta d_n)$ jest sumowalny tą metodą do $\alpha A+\beta D$.
\item Dla każdego całkowitego $m\ge0$ szereg $\sum_{n\ge0}a_n$ w $K$ jest sumowalny tą metodą wtedy i tylko wtedy, gdy szereg $\sum_{n\ge0}a_{n+m}$ jest sumowalny tą metodą. Jeśli suma pierwszego wynosi $A$, suma drugiego wynosi $A-\sum_{j=0}^{m-1}a_j$, przy konwencji sumy pustej równej $0$.
\end{enumerate}
\end{twierdzenie}
\begin{proof}
(i) Dla metody Cesàro stosujemy liniowość skończonych sum i granic. Dla metody Abela przy ustalonym $r<1$ stosujemy liniowość zbieżnych szeregów, a następnie liniowość granicy przy $r\to1^-$.

(ii) Niech $m\ge1$, $h=\sum_{j=0}^{m-1}a_j$ oraz $s_n=\sum_{j=0}^n a_j$. Sumy częściowe ogona to $s_{m+k}-h$. Oznaczmy przez $c_n$ i $d_n$ średnie sum częściowych szeregu pierwotnego i jego ogona. Bezpośrednio ze skończonych sum mamy
\[
 d_n=\frac{n+m+1}{n+1}c_{n+m}
 -\frac{\sum_{j=0}^{m-1}s_j}{n+1}-h.
\]
Współczynnik przy $c_{n+m}$ dąży do $1$, drugi składnik do zera. Stąd $c_n\to A$ implikuje $d_n\to A-h$. Rozwiązując równość względem $c_{n+m}$, otrzymujemy odwrotnie, że $d_n\to D$ implikuje $c_{n+m}\to D+h$, a zatem $c_n\to D+h$.

Dla metody Abela przy $0<r<1$ zachodzi
\[
 \sum_{n=0}^\infty a_nr^n
 =\sum_{j=0}^{m-1}a_jr^j+r^m\sum_{n=0}^\infty a_{n+m}r^n,
\]
ilekroć którykolwiek z szeregów nieskończonych jest zbieżny. Istotnie, odjęcie skończonego początku i dzielenie przez niezerowe $r^m$ pokazują równoważność zbieżności. Dla $r=0$ oba szeregi są automatycznie zbieżne. Skończona suma dąży do $h$, a $r^m\to1$, co po przejściu do granicy w obu kierunkach daje żądany wynik. Dla $m=0$ teza jest natychmiastowa.
\end{proof}
\begin{uwaga}
Regularność i liniowość nie pozwalają dowolnie grupować składników szeregu rozbieżnego. Grupowanie $(1-1)+(1-1)+\cdots$ daje zwykłą sumę $0$, mimo że pierwotny szereg ma sumę Cesàro i Abela równą $1/2$. Grupowanie zmienia ciąg składników, a więc również ciąg sum częściowych. W każdym zastosowaniu należy jawnie wskazać metodę sumowania.
\end{uwaga}
